\section{Historical Records Review}

\subsection{Historical Land Use Summary}
The history of the Subject Property was reconstructed using aerial photographs, topographic maps, city directories, and state archival records. Table~\ref{tab:history} provides a chronological summary of the site's historical land use.

\begin{table}[h]
\centering
\caption{Subject Property Historical Land Use Summary}
\begin{tabularx}{\textwidth}{l X p{3cm}}
\toprule
\textbf{Period} & \textbf{Site Use and Description} & \textbf{Information Source} \\
\midrule
1940 -- 1958 & Agricultural use (orchard and general farming). No structures. & Historical Aerials \\
1958 -- 1974 & Light machining and parts manufacturing (Apex Machine Shop). & City Directories \\
1974 -- 1998 & Electroplating and metal finishing (Synapse Plating Corp.). & State Archive / Permit \\
1998 -- Present & Light electronics assembly and warehousing (Vertex Assemblies). & Interviews / Permits \\
\bottomrule
\end{tabularx}
\label{tab:history}
\end{table}

\subsection{Review of Historical Aerial Photographs}
Historical aerial photographs from 1951, 1968, 1979, 1992, 2005, and 2018 were reviewed to identify physical changes and land use features on and surrounding the Subject Property.
\begin{itemize}
  \item \textbf{1951 Photograph:} The Subject Property and all adjoining properties are under agricultural use. An orchard is visible on the Subject Property.
  \item \textbf{1968 Photograph:} A small commercial/industrial building is visible in the center of the Subject Property. Industrial Parkway has been constructed.
  \item \textbf{1979 Photograph:} The original building has been expanded to its current footprint. A rail spur is visible along the southern boundary. An adjoining commercial building is visible to the north (Nexa Cleaners).
  \item \textbf{1992 Photograph:} The Subject Property appears in its current configuration. A parking lot is located on the west side, and a drum storage pad is visible on the north side.
  \item \textbf{2018 Photograph:} No major structural changes. The area north of the parking lot shows ground disturbance consistent with the UST removal operations.
\end{itemize}

\subsection{Regulatory Database Search}
An environmental database search report was obtained from a commercial data provider on July 28, 2026. The database search compiled records from federal, state, and local regulatory agencies. The results of the search are summarized in Table~\ref{tab:databases}.

\begin{table}[h]
\centering
\caption{Summary of Environmental Database Search Results}
\begin{tabularx}{\textwidth}{X c c c p{4.5cm}}
\toprule
\textbf{Database} & \textbf{Search Distance} & \textbf{Subject Property} & \textbf{Surrounding Sites} & \textbf{Identified Listings} \\
\midrule
Federal NPL & 1.0 mile & 0 & 0 & None \\
Federal CERCLIS / SEMS & 0.5 mile & 0 & 1 & Nexa Cleaners (0.1 mi N) \\
RCRA Hazardous Waste Gen. & Target \& Adjoining & 1 & 2 & Vertex (LQG), Nexa (SQG) \\
State/Tribal LUST & 0.5 mile & 0 & 1 & Meridian Fuel (0.25 mi E) \\
Registered UST/AST & Target \& Adjoining & 1 & 1 & Subject Property (Closed UST) \\
State Spills Database & 0.25 mile & 0 & 1 & Nexa Cleaners (PCE spill, 2011) \\
\bottomrule
\end{tabularx}
\label{tab:databases}
\end{table}

\subsection{Discussion of Surrounding Properties}
The database review identified two surrounding properties with historical or active environmental listings:
\begin{enumerate}
  \item \textbf{Nexa Cleaners (1001 Industrial Parkway):} Located approximately 500 feet north (upgradient hydraulic position). Listed in the State Spills database for a release of PCE in 2011. Remediation is active, and groundwater monitoring has shown a plume migrating south-southeast toward the Subject Property.
  \item \textbf{Meridian Fuel (1050 Industrial Parkway):} Located 1,300 feet east (downgradient). A LUST case was opened in 2005 due to gasoline leakage. The case was closed by the state in 2015 after successful soil excavation. Due to its downgradient position, this site is not considered a migration threat.
\end{enumerate}

